\section{Topic}

\subsection{Introduction to the topic}

Human trafficking is one of the most pervasive transnational crimes in the world, and conflict and post-conflict zones have become some of the most vulnerable environments in which it operates. Armed conflict weakens governing institutions, displaces populations on a mass scale, and disrupts economies, creating conditions in which criminal networks can act with relative impunity. Refugees and internally displaced persons are especially vulnerable, as they often lack documentation, stable housing, employment and access to legal protection, conditions traffickers exploit through deceptive job offers, fraudulent transportation arrangements and false promises of safety.\footnote{United Nations Office on Drugs and Crime, Global Report on Trafficking in Persons 2024 (Vienna: UNODC, 2024) \url{https://www.unodc.org/unodc/en/data-and-analysis/glotip.html}}

A distinguishing feature of trafficking in conflict settings, compared to trafficking in stable states, is the direct involvement of armed actors. Non-state armed groups and terrorist organisations have used trafficking not only as a source of revenue but as an instrument of warfare, employing forced marriage, sexual slavery, forced labour and the recruitment of child soldiers to sustain operations and exert control over civilian populations. This reality was formally recognised by the Security Council in 2016 through Resolution 2331, the first resolution to explicitly link trafficking in persons to armed conflict and terrorism.\footnote{United Nations Security Council, Resolution 2331 (2016), S/RES/2331 (20 December 2016) \url{https://main.un.org/securitycouncil/en/s/res/2331-\%282016\%29}}

Because trafficking in conflict zones sits at the intersection of humanitarian protection, criminal justice and international peace and security, addressing it requires coordination between UNODC, humanitarian actors such as UNHCR and IOM, peacekeeping missions, and the national authorities of source, transit and destination states, the same set of actors examined throughout this guide.

\subsection{Historical Background}

International concern over trafficking dates back over a century, though its legal definition and its connection to armed conflict have evolved considerably over time.

\begin{itemize}
  \item \textbf{Early international efforts (1904–1949)} International concern over trafficking, then framed as the "trafficking of women and children" or "white slave traffic," began with the 1904 International Agreement for the Suppression of the White Slave Traffic, followed by further conventions in 1910, 1921 and 1933. These efforts culminated in the 1949 UN Convention for the Suppression of the Traffic in Persons and of the Exploitation of the Prostitution of Others, the first comprehensive international instrument on the issue, though it lacked a clear, universally accepted definition of trafficking.
  \item \textbf{Cold War and regional conflicts (1950s–1990s)}. Trafficking connected to conflict was documented but rarely addressed as a distinct international priority during this period. The wars in the Balkans during the 1990s, however, drew significant international attention to the link between armed conflict, mass displacement and sexual exploitation, as post-conflict Bosnia and Kosovo became notable destination and transit points for trafficking networks exploiting the presence of international peacekeeping and reconstruction personnel.
  \item \textbf{Codiffication of a modern legal deffinition (2000)}. The adoption of the Palermo Protocol in 2000 marked a turning point, establishing the first internationally accepted legal definition of trafficking in persons and obliging states to criminalise it, protect victims and cooperate across borders. This Protocol remains the primary legal framework guiding the Committee's work today.
  \item \textbf{Trafficking by armed and terrorist groups (2014–2017)}. The rise of ISIS in Iraq and Syria brought global attention to the systematic use of trafficking by non-state armed actors, most notably the enslavement of Yazidi women and girls beginning in 2014. This period prompted the Security Council to adopt Resolution 2331 (2016) and Resolution 2388 (2017), formally recognising trafficking in conflict as a threat to international peace and security.
\end{itemize}

\subsubsection{Displacement-driven trafficking in recent conflicts}

\textbf{(2022–present)}. The wars in Ukraine, Sudan and other ongoing conflicts have demonstrated how rapidly trafficking risks emerge during large-scale displacement crises, with international organisations warning of traffickers exploiting refugees through fraudulent employment schemes, online recruitment and exploitation within displacement settlements.\footnote{United Nations Office on Drugs and Crime, Global Report on Trafficking in Persons 2024 (Vienna: UNODC, 2024) \url{https://www.unodc.org/unodc/en/data-and-analysis/glotip.html}}

\subsection{Past Actions}

The international community has taken numerous actions over the past two decades to combat human trafficking in relation to armed conflict and transnational criminal networks. These efforts have focused on strengthening international legal frameworks, improving cooperation between states, protecting victims and addressing trafficking networks linked to conflict and terrorism.

One of the most significant international agreements addressing human trafficking is the Protocol to Prevent, Suppress and Punish Trafficking in Persons Especially Women and Children, a.k.a the Palermo Protocol. Adopted by the UN General Assembly in 2000 as part of the United Nations Convention against Transnational Organised Crime, the Palermo Protocol established the first internationally accepted definition of human trafficking.\footnote{United Nations Office on Drugs and Crime. (2000). Protocol to prevent, suppress and punish trafficking in persons, especially women and children, supplementing the United Nations Convention against Transnational Organized Crime. United Nations. \url{https://www.unodc.org/unodc/en/human-trafficking/protocol.html}} It obliged states to criminalise trafficking, protect victims, strengthen border controls and work to cooperate internationally to combat human trafficking networks.\footnote{United Nations Office on Drugs and Crime. (2000). Protocol to prevent, suppress and punish trafficking in persons, especially women and children, supplementing the United Nations Convention against Transnational Organized Crime. United Nations. \url{https://www.unodc.org/unodc/en/human-trafficking/protocol.html}} Today, the Palermo Protocol is one of the primary international legal frameworks guiding anti-trafficking policies.\footnote{United Nations Office on Drugs and Crime. (2000). Protocol to prevent, suppress and punish trafficking in persons, especially women and children, supplementing the United Nations Convention against Transnational Organized Crime. United Nations. \url{https://www.unodc.org/unodc/en/human-trafficking/protocol.html}}

In 2010, the United Nations General Assembly adopted the Global Plan of Action to Combat Trafficking in Persons.\footnote{United Nations General Assembly. (2010). United Nations Global Plan of Action to Combat Trafficking in Persons (A/RES/64/293, 12 August 2010). United Nations. \url{https://digitallibrary.un.org/record/687201}} This plan encouraged member states to coordinate efforts regarding prevention, prosecution and victim protection while also emphasising the importance of addressing the social and economic conditions that contribute to trafficking.\footnote{United Nations General Assembly. (2010). United Nations Global Plan of Action to Combat Trafficking in Persons (A/RES/64/293, 12 August 2010). United Nations. \url{https://digitallibrary.un.org/record/687201}} The initiative also established the United Nations Voluntary Trust Fund for Victims of Trafficking in Persons, which supports organisations providing assistance to survivors around the world.\footnote{United Nations General Assembly. (2010). United Nations Global Plan of Action to Combat Trafficking in Persons (A/RES/64/293, 12 August 2010). United Nations. \url{https://digitallibrary.un.org/record/687201}}

The United Nations Security Council has increasingly recognised trafficking in persons in conflict and post-conflict situations as a threat to international peace and security. In 2016, the Security Council adopted Resolution 2331, the first resolution specifically linking human trafficking to armed conflict and terrorism.\footnote{United Nations Security Council. (2016). Resolution 2331 (S/RES/2331, 20 December 2016). United Nations. \url{https://main.un.org/securitycouncil/en/s/res/2331-\%282016\%29}} The resolution condemned trafficking by terrorist groups such as ISIS and Boko Haram, recognising that these organisations used trafficking for sexual slavery, forced labour, recruitment and financing operations.\footnote{United Nations Security Council. (2016). Resolution 2331 (S/RES/2331, 20 December 2016). United Nations. \url{https://main.un.org/securitycouncil/en/s/res/2331-\%282016\%29}} Resolution 2331 called upon member states to improve information-sharing, strengthen border security, investigate trafficking linked to terrorism and protect vulnerable populations affected by conflict.\footnote{United Nations Security Council. (2016). Resolution 2331 (S/RES/2331, 20 December 2016). United Nations. \url{https://main.un.org/securitycouncil/en/s/res/2331-\%282016\%29}} In 2017, the Security Council adopted Resolution 2388, which expanded international efforts to address trafficking in conflict zones.\footnote{United Nations Security Council. (2017). Resolution 2388 (S/RES/2388, 21 November 2017). United Nations. \url{https://main.un.org/securitycouncil/en/content/sres2388-2017}} The resolution urged member states to improve victim identification, prosecute traffickers more effectively and strengthen cooperation between governments, humanitarian agencies and law enforcement organisations and institutions. It also emphasized the importance of specifically supporting women and children affected by armed conflict and displacement.\footnote{United Nations Security Council. (2017). Resolution 2388 (S/RES/2388, 21 November 2017). United Nations. \url{https://main.un.org/securitycouncil/en/content/sres2388-2017}}

The UNODC has played a central role in supporting member states through technical and training assistance and research initiatives.\footnote{United Nations Office on Drugs and Crime. (n.d.). Human trafficking and migrant smuggling. UNODC. \url{https://www.unodc.org/unodc/human-trafficking/}} UNODC regularly publishes the Global Report on Trafficking in Persons, which provides data and analyses on trafficking trends worldwide.\footnote{United Nations Office on Drugs and Crime. (2024). Global Report on Trafficking in Persons 2024. United Nations. \url{https://www.unodc.org/documents/data-and-analysis/glotip/2024/GLOTIP2024_BOOK.pdf}} The organisation also assists governments in drafting anti-trafficking legislation,  training  police  and  judicial  officials  to  help  improve international legal cooperation.\footnote{United Nations Office on Drugs and Crime. (n.d.). Human trafficking and migrant smuggling. UNODC. \url{https://www.unodc.org/unodc/human-trafficking/}} The IOM has implemented numerous anti-trafficking initiatives in migration passages and conflict-affected regions.\footnote{International Organization for Migration. (2024). World Migration Report 2024 (M. McAuliffe \& L.A. Oucho, Eds.). IOM. \url{https://publications.iom.int/books/world-migration-report-2024}} IOM works closely with governments and humanitarian agencies to identify trafficking victims and provide them with rehabilitation services, working to improve migrant protection systems.\footnote{International Organization for Migration. (2024). World Migration Report 2024 (M. McAuliffe \& L.A. Oucho, Eds.). IOM. \url{https://publications.iom.int/books/world-migration-report-2024}} In conflict settings, IOM has established monitoring systems to track displacement trends and identify populations at heightened risk of exploitation.\footnote{International Organization for Migration. (n.d.). Displacement Tracking Matrix: National displacement reports. IOM. \url{https://dtm.iom.int/report-product-series/national-displacement-report}}

Regional organisations have also developed anti-trafficking frameworks. The EU adopted the EU Anti-Trafficking Directive in 2011, strengthening protections for victims and improving coordination among member states.\footnote{Directive 2011/36/EU of the European Parliament and of the Council of 5 April 2011 on preventing and combating trafficking in human beings and protecting its victims, and replacing Council Framework Decision 2002/629/JHA, OJ L 101, 15.4.2011. \url{https://eur-lex.europa.eu/eli/dir/2011/36/oj/eng}} The African Union has supported regional cooperation initiatives focused on trafficking prevention and migration management, while ASEAN has adopted regional conventions addressing trafficking in Southeast Asia.

Human trafficking has also been incorporated into broader international development goals. In 2015, the United Nations included anti-trafficking targets within the Sustainable Development Goals (SDGs).\footnote{United Nations General Assembly. (2015). Transforming our world: The 2030 Agenda for Sustainable Development (A/RES/70/1). United Nations. \url{https://sdgs.un.org/2030agenda}} SDG 5 calls for the elimination of violence against women and girls, including trafficking and sexual exploitation.\footnote{United Nations General Assembly. (2015). Transforming our world: The 2030 Agenda for Sustainable Development (A/RES/70/1). United Nations. \url{https://sdgs.un.org/2030agenda}} SDG 8.7 urges states to eradicate forced labour, modern slavery and child trafficking, while SDG 16 focuses on reducing exploitation of and abuse against children.\footnote{United Nations. (n.d.). Goal 8: Decent work and economic growth — Target 8.7. United Nations Department of Economic and Social Affairs. \url{https://sdgs.un.org/goals/goal8}; see also United Nations. (n.d.). Goal 16: Peace, justice and strong institutions — Target 16.2. \url{https://sdgs.un.org/goals/goal16}} Peacekeeping operations and humanitarian missions have increasingly integrated anti-trafficking measures into their mandates. UN agencies and NGOs now provide training for humanitarian workers and peacekeepers to recognize indicators of trafficking and protect vulnerable populations in refugee camps and post-conflict communities.\footnote{United Nations Office on Drugs and Crime. (2024). Global Report on Trafficking in Persons 2024. United Nations. \url{https://www.unodc.org/documents/data-and-analysis/glotip/2024/GLOTIP2024_BOOK.pdf}} Some missions also cooperate with local authorities to strengthen law enforcement capacity and victim support systems.

Despite these efforts, major challenges remain. Many anti-trafficking initiatives suffer from insufficient funding, inconsistent implementation and limited inter-state coordination.\footnote{United Nations Office on Drugs and Crime. (2024). Global Report on Trafficking in Persons 2024. United Nations. \url{https://www.unodc.org/documents/data-and-analysis/glotip/2024/GLOTIP2024_BOOK.pdf}} Weak governance in conflict zones continues to hinder prosecutions and victim protection efforts. In addition, rapid technological advancements have allowed traffickers to expand online recruitment methods faster than many governments can regulate them.\footnote{United Nations Office on Drugs and Crime. (2024). Global Report on Trafficking in Persons 2024. United Nations. \url{https://www.unodc.org/documents/data-and-analysis/glotip/2024/GLOTIP2024_BOOK.pdf}} Nevertheless, past actions by the UN, regional organisations and individual states have significantly increased international awareness of trafficking in conflict zones. These efforts have helped to establish important legal frameworks, improve international cooperation while recognising trafficking as both a serious humanitarian and security issue.

\subsection{Current Situation and Statistics}

Human trafficking is one of the most prevalent transnational crimes in the world, and conflict and post-conflict zones have become particularly vulnerable to trafficking networks.\footnote{United Nations Office on Drugs and Crime. (2024). Global Report on Trafficking in Persons 2024. United Nations. \url{https://www.unodc.org/documents/data-and-analysis/glotip/2024/GLOTIP2024_BOOK.pdf}} Armed conflict weakens governmental institutions, displaces and destabilises populations and disrupts economies, creating conditions whereby criminal organisations can operate with relatively little consequences.\footnote{United Nations Office on Drugs and Crime. (2024). Global Report on Trafficking in Persons 2024. United Nations. \url{https://www.unodc.org/documents/data-and-analysis/glotip/2024/GLOTIP2024_BOOK.pdf}} In recent years, the international community has observed a considerable increase in trafficking linked to wars and humanitarian crises.\footnote{United Nations Office on Drugs and Crime. (2024). Global Report on Trafficking in Persons 2024. United Nations. \url{https://www.unodc.org/documents/data-and-analysis/glotip/2024/GLOTIP2024_BOOK.pdf}}

According to the 2024 Global Report on Trafficking in Persons published by UNODC, global trafficking cases rose significantly after the COVID-19 pandemic, with detected victims increasing by around 25\% compared to 2019 statistics.\footnote{United Nations Office on Drugs and Crime. (2024, December 16). UNODC global human trafficking report: Detected victims up 25 per cent as more children are exploited and forced labour cases spike [Press release]. \url{https://www.unodc.org/unodc/en/press/releases/2024/December/unodc-global-human-trafficking-report_-detecte} d-victims-up-25-per-cent-as-more-children-are-exploited-and-forced-labour-cases-spike.html} The report identifies some major factors that are contributing to the vulnerability and exploitation of people, including conflict, poverty, climate disasters and economic instability.\footnote{United Nations Office on Drugs and Crime. (2024). Global Report on Trafficking in Persons 2024. United Nations. \url{https://www.unodc.org/documents/data-and-analysis/glotip/2024/GLOTIP2024_BOOK.pdf}} Additionally, women and children continue to represent the majority of trafficking victims, particularly in conflict-affected regions where their legal protections and social support services are limited.\footnote{United Nations Office on Drugs and Crime. (2024). Global Report on Trafficking in Persons 2024. United Nations. \url{https://www.unodc.org/documents/data-and-analysis/glotip/2024/GLOTIP2024_BOOK.pdf}}

Conflict zones create ideal conditions for traffickers because state institutions are often weakened or entirely nonfunctional.\footnote{United Nations Office on Drugs and Crime. (2024). Global Report on Trafficking in Persons 2024. United Nations. \url{https://www.unodc.org/documents/data-and-analysis/glotip/2024/GLOTIP2024_BOOK.pdf}} Judicial systems and security officials may become ineffective during and post-conflicts in countries where violence, corruption or political instability has shaken the foundations of governments and their ability to protect their citizens from predators. As a result, traffickers can more easily move their victims across borders, exploiting displaced persons with limited risk of prosecution.\footnote{United Nations Office on Drugs and Crime. (2024). Global Report on Trafficking in Persons 2024. United Nations. \url{https://www.unodc.org/documents/data-and-analysis/glotip/2024/GLOTIP2024_BOOK.pdf}} Refugees and internally displaced persons (IDPs) are especially vulnerable because they often lack proper documentation, stable employment, housing, education and access to legal protections.\footnote{United Nations High Commissioner for Refugees. (2024). Global Trends: Forced Displacement in 2023. UNHCR. \url{https://www.unhcr.org/sites/default/files/2024-06/global-trends-report-2023.pdf}}

The ongoing conflicts in Syria, Sudan, Ukraine, Afghanistan, Libya and the Democratic Republic of the Congo illustrate how war and instability intensify trafficking risks.\footnote{United Nations Office on Drugs and Crime. (2024). Global Report on Trafficking in Persons 2024. United Nations. \url{https://www.unodc.org/documents/data-and-analysis/glotip/2024/GLOTIP2024_BOOK.pdf}} Millions of displaced individuals fleeing these conflicts have become targets for forced labour, sexual exploitation, child soldier recruitment, forced marriage and organ trafficking.\footnote{United Nations Office on Drugs and Crime. (2024). Global Report on Trafficking in Persons 2024. United Nations. \url{https://www.unodc.org/documents/data-and-analysis/glotip/2024/GLOTIP2024_BOOK.pdf}} In refugee camps and settlements, traffickers often exploit shortages of employment and humanitarian assistance by offering deceptive job opportunities or transportation arrangements.\footnote{United Nations Office on Drugs and Crime. (2024). Global Report on Trafficking in Persons 2024. United Nations. \url{https://www.unodc.org/documents/data-and-analysis/glotip/2024/GLOTIP2024_BOOK.pdf}} The war in Ukraine has demonstrated how rapidly trafficking risks can emerge during large-scale displacement crises.

Since 2022, millions of refugees (primarily women and children) have crossed into neighbouring European states.\footnote{United Nations High Commissioner for Refugees. (2024). Global Trends: Forced Displacement in 2023. UNHCR. \url{https://www.unhcr.org/sites/default/files/2024-06/global-trends-report-2023.pdf}} International organisations including the IOM and UNHCR warned early in the conflict that traffickers were targeting vulnerable refugees through fraudulent employment and transportation schemes as well as recruitment tactics via the internet.\footnote{UNHCR. (2022, April 12). Statement on risks of trafficking and exploitation facing refugees from Ukraine [Press release]. \url{https://www.unhcr.org/us/news/press-releases/statement-risks-trafficking-and-exploitation-facing-refugees-ukrai} ne-attributed; see also UNODC. (2022, March). Targeted by traffickers — Ukrainian refugees at high risk of exploitation [Press release]. \url{https://www.unodc.org/unodc/press/releases/2022/March/targeted-by-traffickers---ukrainian-refugees-at-high-ris} k-of-exploitation.html} Similar patterns have appeared in conflicts across the Middle East and Africa, particularly in regions with weak border controls and high levels of displacement.\footnote{International Organization for Migration. (2024). World Migration Report 2024 (M. McAuliffe \& L.A. Oucho, Eds.). IOM. \url{https://publications.iom.int/books/world-migration-report-2024}}

Armed groups and terrorist organisations have also been identified as using trafficking as both a weapon and a source of revenue. Groups such as ISIS exploited women and children through forced marriages and recruitment into armed conflict.\footnote{United Nations Human Rights Council. (2016). "They came to destroy": ISIS crimes against the Yazidis (A/HRC/32/CRP.2, 15 June 2016). \url{https://www.ohchr.org/en/hr-bodies/hrc/regular-sessions/session32/res-dec-stat}} In several regions of Africa and the Middle East, extremist organisations continue to finance their operations through human trafficking and forced labour. Criminal networks are becoming increasingly cooperative across borders, making trafficking more difficult for individual states to combat independently.\footnote{United Nations Office on Drugs and Crime. (2024). Global Report on Trafficking in Persons 2024. United Nations. \url{https://www.unodc.org/documents/data-and-analysis/glotip/2024/GLOTIP2024_BOOK.pdf}}

Technology has also become an important tool in trafficking networks. Traffickers now use social media platforms, encrypted messaging applications and online job advertisements to recruit victims.\footnote{United Nations Office on Drugs and Crime. (2024). Global Report on Trafficking in Persons 2024. United Nations. \url{https://www.unodc.org/documents/data-and-analysis/glotip/2024/GLOTIP2024_BOOK.pdf}} Digital platforms allow criminal networks to operate internationally while concealing their identities and locations. The rise of online scams and internet-enabled forced labour has become particularly alarming in recent years, such as traffickers exploiting displaced populations through false employment schemes.\footnote{United Nations Office on Drugs and Crime. (2024). Global Report on Trafficking in Persons 2024. United Nations. \url{https://www.unodc.org/documents/data-and-analysis/glotip/2024/GLOTIP2024_BOOK.pdf}} Despite increased global awareness, many states continue to face serious obstacles in responding effectively to trafficking in conflict and post-conflict zones. In many post-conflict countries, governments lack the financial resources and institutional capacity necessary to identify victims, prosecute traffickers and provide victims with rehabilitation services.\footnote{United Nations Office on Drugs and Crime. (2024). Global Report on Trafficking in Persons 2024. United Nations. \url{https://www.unodc.org/documents/data-and-analysis/glotip/2024/GLOTIP2024_BOOK.pdf}} Corruption within law enforcement and border agencies further undermines possible anti-trafficking efforts, and many victims fear reporting crimes because of stigma and distrust of authorities.

Humanitarian agencies have also struggled with funding shortages that affect anti-trafficking programs. In 2025, the IOM warned that major cuts to humanitarian funding threatened the programs assisting vulnerable migrants and trafficking victims worldwide.\footnote{International Organization for Migration. (2025, March 18). Update on IOM operations amid budget cuts [Statement]. \url{https://www.iom.int/news/update-iom-operations-amid-budget-cuts}} Reduced funding risks a further weakening of protection systems within refugee camps, conflict zones and popular passages used by displaced persons and migrants, leaving these already vulnerable populations at even greater risk of exploitation.\footnote{International Organization for Migration. (2025, March 18). Update on IOM operations amid budget cuts [Statement]. \url{https://www.iom.int/news/update-iom-operations-amid-budget-cuts}}

International law provides important frameworks for combating trafficking, but its implementation is inconsistent. While many states have ratified the Palermo Protocol, enforcement mechanisms and victim support systems differ greatly from country to country.\footnote{United Nations Office on Drugs and Crime. (2024). Global Report on Trafficking in Persons 2024. United Nations. \url{https://www.unodc.org/documents/data-and-analysis/glotip/2024/GLOTIP2024_BOOK.pdf}} Cross-border cooperation is limited in many conflict-affected regions, allowing trafficking networks to more easily exploit legal loopholes and jurisdictional gaps. As conflicts continue to expand and displacement levels are historically high, human trafficking is increasingly recognised not only as a human rights issue, but also a threat to international peace and security.\footnote{United Nations Security Council. (2016). Resolution 2331 (S/RES/2331, 20 December 2016). United Nations. \url{https://main.un.org/securitycouncil/en/s/res/2331-\%282016\%29}} The challenge facing the international community is to strengthen prevention mechanisms, improve international coordination and support and address the root causes that allow trafficking to flourish in conflict and post-conflict environments.

\subsection{Problem Statement}

Human trafficking in conflict and post-conflict zones represents a serious threat to international peace, security and human rights. The issue is driven by several factors:

\begin{itemize}
  \item \textbf{Sexual exploitation} (e.g., women and girls forced into prostitution by armed groups or criminal networks)
  \item \textbf{Forced labour slavery} (including children)
\end{itemize}

(e.g., children forced to work in mines, agriculture, or domestic service and etc.)

\subsubsection{Recruitment and use of child soldiers}

(e.g., children coerced into fighting, spying, or carrying supplies for armed groups)

\subsubsection{Forced marriage}

(e.g., girls abducted and forced to marry members of armed groups)

\subsubsection{Organ trafficking}

e.g., vulnerable people deceived or coerced into having organs removed for illegal sale)

\subsubsection{Human trafficking among refugees and IDPs}

(e.g., displaced people lured by false promises of jobs or safe passage and then exploited)

\subsubsection{Mass displacement}

(e.g., families fleeing conflict become separated, making individuals easier targets for traffickers)

\subsubsection{Poverty and economic instability}

(e.g., people accepting fraudulent job offers due to lack of income)

\subsubsection{Limited resources for victims}

(e.g., shortages of shelters, legal aid, or healthcare preventing victims from receiving protection and support)

These challenges make it difficult for states to identify victims and prevent trafficking in regions already destabilized by conflict.

\subsection{Bloc Positions}

While many countries agree on major points regarding combatting human trafficking in conflict and post-conflict zones, there are still multiple issues of contention, as different states prioritize different parts of the problem. Consequently, all members of our committee could be divided into several blocs based on their perspectives on the problem at hand.

\subsubsection{Conflict-Affected Source States}

Countries: Afghanistan, Central African Republic, Chad, Democratic Republic of Congo, Iraq, Libya, Mali, Myanmar, Somalia, South Sudan, Sudan, Syrian Arab Republic, Ukraine, Yemen

For these countries that experience armed conflict or emerge from prolonged instability, human trafficking is often deeply intertwined with the realities of war. The breakdown of governance structures, weakened law enforcement capacities, widespread displacement, and economic collapse create nearly perfect conditions in which trafficking networks can thrive.\footnote{United Nations Office on Drugs and Crime, Global Report on Trafficking in Persons 2024 (Vienna: UNODC, 2024) \url{https://www.unodc.org/unodc/en/data-and-analysis/glotip.html}} Trafficking is not merely a by-product of conflict, it is a tool used by armed actors to sustain military operations and exert control over civilian populations.

For instance, in the DRC, armed groups have long relied on forced labor in mining regions, while children have been recruited as combatants, porters, and informants.\footnote{United Nations Children's Fund (UNICEF), "Children recruited by armed forces or armed groups" (UNICEF, updated 2024) \url{https://www.unicef.org/protection/children-recruited-by-armed-forces}} Similar trends have been documented in South Sudan, Somalia, and Myanmar, where armed factions have recruited children into military service.\footnote{Office of the Special Representative of the Secretary-General for Children and Armed Conflict, Annual Report of the Secretary-General on Children and Armed Conflict (UN Doc A/78/842–S/2024/384, June 2024). Available at \url{https://childrenandarmedconflict.un.org}} In Iraq and Syria, thousands of Yazidi women and girls were subjected to systematic sexual slavery by ISIS.\footnote{United Nations Human Rights Council, "They came to destroy": ISIS Crimes Against the Yazidis, Report of the Independent International Commission of Inquiry on the Syrian Arab Republic (UN Doc A/HRC/32/CRP.2, 15 June 2016). This is the landmark 2016 report finding that ISIS committed genocide against the Yazidis. Available at \url{https://www.ohchr.org/en/hr-bodies/hrc/regular-sessions/session32/res-dec-stat}} More recently, the wars in Sudan and Ukraine have generated large-scale displacement, increasing the vulnerability of refugees and IDPs to traffickers operating along migration routes.\footnote{United Nations High Commissioner for Refugees, Global Trends: Forced Displacement in 2023 (Geneva: UNHCR, 2024) \url{https://www.unhcr.org/global-trends}} \footnote{International Organization for Migration, World Migration Report 2024 (Geneva: IOM, 2024) \url{https://publications.iom.int/books/world-migration-report-2024}}

\subsubsection{Refugee-Hosting and Frontline States}

Countries: Jordan, Lebanon, Türkiye, Egypt, Ethiopia, Chad, Pakistan

These countries host significant populations of refugees and IDPs fleeing nearby conflicts. They often function as both destination and transit countries for vulnerable populations. These states are particularly concerned with the trafficking risks faced by refugees, especially women, children, and unaccompanied minors, who may be targeted by traffickers during migration or within refugee camps.\footnote{United Nations High Commissioner for Refugees, Global Trends: Forced Displacement in 2023 (2024) \url{https://www.unhcr.org/sites/default/files/2024-06/global-trends-report-2023.pdf}} Their primary priorities include greater international burden-sharing, increased humanitarian funding, stronger refugee registration systems, and enhanced cooperation with international organisations. Refugee-hosting and frontline states stress that the international community must do more to support countries that bear the social and economic costs of displacement.\footnote{United Nations High Commissioner for Refugees, Global Trends: Forced Displacement in 2023 (2024) \url{https://www.unhcr.org/sites/default/files/2024-06/global-trends-report-2023.pdf}}

\subsubsection{Destination Countries and Human Rights Advocates}

Countries: Canada, France, Germany, Norway, Sweden, United Kingdom, United States of America

The countries of this bloc comprise major destination countries for migrants and refugees and generally adopt a human rights-centered approach to combating trafficking. The states emphasize the protection of victims, particularly survivors of sexual exploitation, forced labor, and gender-based violence, while advocating for stronger international accountability mechanisms.\footnote{United Nations Office of Drugs and Crime, Global Report on Trafficking in Persons 2024 (2024). \url{https://www.unodc.org/documents/data-and-analysis/glotip/2024/GLOTIP2024_BOOK.pdf}} Members of this group support such initiatives as international monitoring frameworks, sanctions against trafficking networks, enhanced prosecution of traffickers, and survivor-centered rehabilitation programs. They also push for stronger legal language regarding the responsibilities of governments and armed groups involved in trafficking-related abuses.\footnote{International Labour Organization, Walk Free, and International Organization for Migration, Global Estimates of Modern Slavery: Forced Labour and Forced Marriage (Geneva: ILO, 2022) \url{https://www.ilo.org/publications/global-estimates-modern-slavery-forced-labour-and-forced-marria} ge-1} Their stance and claims tend to clash with the ones of another bloc that prioritizes sovereignty over international oversight.

\subsubsection{Sovereignty-Oriented and Security-Focused States}

Countries: China, Russia, India, Pakistan, Egypt, Myanmar

These countries generally support international cooperation against trafficking. However, they maintain that responses should respect state sovereignty and the principle of non-interference in domestic affairs. Rather than emphasizing international monitoring mechanisms, the member states of this group prioritize strengthening national institutions, improving border security and enhancing law enforcement cooperation, and combating terrorism and TOC.\footnote{United Nations Office of Drugs and Crime, Global Report on Trafficking in Persons 2024 (2024). \url{https://www.unodc.org/documents/data-and-analysis/glotip/2024/GLOTIP2024_BOOK.pdf}} These states are often cautious about solutions and projects that create external accountability structures or impose country-specific scrutiny. Instead, they advocate for technical assistance, information sharing, and nationally led anti-trafficking initiatives.\footnote{United Nations General Assembly, United Nations Convention against Transnational Organized Crime (res 55/25, 15 November 2000). The UNTOC framework establishes principles of technical assistance and state-led implementation. Available at \url{https://www.unodc.org/unodc/en/organized-crime/intro/UNTOC.html}}

\subsubsection{Transit States and Organized Crime-Affected States}

Countries: Brazil, Colombia, Mexico, Nigeria, Philippines, United Arab Emirates, Libya, Türkiye

These countries occupy important positions along major migration and trafficking routes and often confront significant challenges posed by TOC networks. They serve as countries of origin, transit, or destination for trafficking victims. The states of this bloc focus mainly on dismantling criminal organizations involved in migrant smuggling, forced labor, and sexual exploitation.\footnote{UNODC, Global Report on Trafficking in Persons 2024 (2024).} The countries of this group support measures such as cross-border intelligence sharing, financial tracking of trafficking profits, enhanced border management, and stronger cooperation between law enforcement agencies.\footnote{International Organization for Migration, World Migration Report 2024 (Geneva: IOM, 2024) \url{https://publications.iom.int/books/world-migration-report-2024}} Many members also advocate for addressing the economic factors that make individuals vulnerable to exploitation.

\subsubsection{Development-Oriented Emerging Powers}

Countries: Brazil, India, China, South Africa, Mexico, Nigeria

Although the approaches of these committee countries vary, these states all emphasize the relationship between trafficking and broader socioeconomic challenges. They argue that conflict-related trafficking cannot be addressed exclusively through criminal justice measures and that long-term solutions require tackling poverty, unemployment, inequality, weak institutions, and underdevelopment.\footnote{International Labour Organization, Walk Free, and International Organization for Migration, Global Estimates of Modern Slavery: Forced Labour and Forced Marriage (Geneva: ILO, 2022) \url{https://www.ilo.org/publications/global-estimates-modern-slavery-forced-labour-and-forced-marria} ge-1} Consequently, they support initiatives focused on economic recovery, education, employment opportunities, and sustainable development as tools for reducing vulnerability to trafficking in post-conflict societies.

\subsubsection{General State of Affairs}

The biggest divide among countries emerges between human rights-focused  destination  states  and  sovereignty-oriented security-focused states. While the former advocate for stronger international accountability and victim protection mechanisms, the latter emphasize national control over implementation. The second divide arises between refugee-hosting countries, which seek greater international assistance, and wealthier donor states, which may prioritize accountability and governance reforms.

Finally, there exists a debate whether the international community should focus primarily on law enforcement and prosecution or on addressing the root causes of trafficking, including conflict, displacement, and economic instability.\footnote{United Nations Office of Drugs and Crime, Global Report on Trafficking in Persons 2024 (2024). \url{https://www.unodc.org/documents/data-and-analysis/glotip/2024/GLOTIP2024_BOOK.pdf}}

\subsection{Key Stakeholders}

While states remain the primary actors responsible for preventing and combating human trafficking, they are not the only stakeholders that influence the issue. Given the complex nature of trafficking in conflict and post-conflict settings, a diverse range of international organisations, civil society actors, as well as non-state armed groups shape both the growth of trafficking networks and the effectiveness of efforts to counter them.

\subsubsection{International Organisations (IOs)}

International  organisations  play  a  central  role  in  coordinating anti-trafficking efforts, particularly in contexts where state institutions have been weakened by conflict, displacement, or political instability. The UNODC serves as the leading international body responsible for supporting the implementation of anti-trafficking legislation, strengthening criminal justice systems, as well as facilitating international cooperation against transnational organized crime.\footnote{United Nations Office of Drugs and Crime, "Human Trafficking and Migrant Smuggling" (UNODC portal) \url{https://www.unodc.org/unodc/en/human-trafficking/index.html}} Alongside UNODC, organisations such as the UNHCR, the IOM, and UNICEF focus on protecting vulnerable populations, including refugees, IDPs, and children affected by armed conflict.\footnote{United Nations High Commissioner for Refugees, Global Trends: Forced Displacement in 2023 (2024) \url{https://www.unhcr.org/sites/default/files/2024-06/global-trends-report-2023.pdf}} \footnote{International Organization for Migration, World Migration Report 2024 (Geneva: IOM, 2024). \url{https://publications.iom.int/books/world-migration-report-2024}}

\subsubsection{Non-Governmental organisations (NGOs)}

Beyond direct assistance, IOs contribute to data collection, capacity-building, humanitarian relief, and post-conflict reconstruction initiatives. Their activities are particularly important in situations where trafficking risks increase due to large-scale displacement, weakened governance, and the breakdown of social protection mechanisms.\footnote{International Organization for Migration, World Migration Report 2024 (Geneva: IOM, 2024). \url{https://publications.iom.int/books/world-migration-report-2024}}

In many conflict-affected regions, NGOs are among the first actors to identify trafficking risks and provide assistance to vulnerable populations. IOs such as Save the Children, World Vision, and the International Rescue Committee (IRC) deliver humanitarian aid, child protection services, legal assistance, and psychosocial support to survivors of exploitation and abuse.\footnote{Save the Children, "What We Do: Child Protection" (Save the Children, 2024) \url{https://www.savethechildren.org.uk/what-we-do}}

In addition to service provision, NGOs play an important advocacy role by documenting human rights violations, raising awareness of trafficking trends, and encouraging governments and international institutions to strengthen protective measures. Their close engagement with local communities often allows them to identify emerging threats more rapidly than state institutions, especially in areas where government authority is limited or absent. Nevertheless, their operations are frequently constrained by funding shortages, security concerns, and restricted access to active conflict zones.\footnote{United Nations Office for the Coordination of Humanitarian Affairs (OCHA), Global Humanitarian Overview 2023 (Geneva: OCHA, 2022) \url{https://www.unocha.org/publications/report/world/global-humanitarian-overview-2023-enaresfr}}

\subsubsection{Non-State Armed Groups and Terrorist organisations}

Non-state armed groups and terrorist organisations represent some of the most significant perpetrators of trafficking in conflict settings. Unlike conventional criminal enterprises, these actors frequently employ trafficking as both a source of revenue and a strategic instrument of warfare. Victims may be subjected to forced labor, sexual slavery, forced marriage, human smuggling, or recruitment into armed service.\footnote{United Nations Human Rights Council, "They came to destroy": ISIS Crimes Against the Yazidis (UN Doc A/HRC/32/CRP.2, 15 June 2016) \url{https://www.ohchr.org/sites/default/files/Documents/HRBodies/HRCouncil/CoISyria/A_HRC_32_} CRP.2\_en.pdf}

The involvement of such groups highlights the close relationship between trafficking, armed conflict, terrorism, and TOC, demonstrating why anti-trafficking efforts in conflict zones frequently intersect with broader peacebuilding and security objectives.

\subsection{Guiding Questions - QARMAs}

\begin{enumerate}
  \item How can vulnerable populations in conflict and post-conflict zones be identified before they become victims of trafficking?
  \item What measures can be implemented within refugee camps and displacement settlements to reduce trafficking risks?
  \item How can the international community address the recruitment and use of child soldiers by armed groups?
  \item What strategies can be used to combat conflict-related sexual exploitation and sexual slavery?
  \item How can states strengthen cross-border cooperation to disrupt trafficking networks operating across conflict zones?
  \item What role should technology play in preventing and detecting human trafficking?
  \item How can victims and survivors of trafficking be effectively rehabilitated and reintegrated into society?
  \item How can post-conflict reconstruction efforts reduce the long-term risk of trafficking?
  \item What actions can be taken to hold traffickers, armed groups, and other perpetrators accountable?
  \item How can international organizations, national governments, and local communities coordinate their efforts to combat human trafficking more effectively?
\end{enumerate}
